% ASSERT_CONVENTION: natural_units=natural, metric_signature=mostly_minus,
%   fourier_convention=physics, coupling_convention=alpha_s,
%   renormalization_scheme=MSbar, gauge_choice=Feynman,
%   spinor_trace=4_in_d4, tr_gamma_munu=0, heat_kernel_a1=E_plus_R_over_6,
%   one_loop_sign=boson_minus_half_Trlog_fermion_plus_half_Trlog,
%   squared_Dirac=Lichnerowicz_E_eq_minus_R_over_4_minus_half_gamma_munu_F,
%   spacetime_block_engine_indices={1,2,3,10}, Vhalf_engine_indices={11,18,19,26},
%   scheme=proper_time_cutoff_vacuum_Casimir_NO_thermodynamics,
%   det_SSOT=ring_lemma_verification.py, harness=bulk_geometry_verification.py+cartan_phaseB_curvature.py
%
% NOTE: pdflatex is NOT available in this execution environment (checked: `which
% pdflatex` -> not found). This source is well-formed (balanced environments, all
% \ref/\eqref resolve to defined \label targets) but was NOT compiled in-env. Compile
% downstream with `pdflatex 81-sakharov-gate1.tex` (no exotic packages required).
\documentclass[11pt]{article}
\usepackage{amsmath,amssymb,amsthm}
\newtheorem{proposition}{Proposition}
\newtheorem{lemma}{Lemma}
\newtheorem{remark}{Remark}
\newcommand{\Tr}{\mathrm{Tr}}
\newcommand{\STr}{\mathrm{STr}}
\newcommand{\h}{\mathfrak{h}}
\newcommand{\Oo}{\mathbb{O}}
\newcommand{\RR}{\mathbb{R}}
\newcommand{\QQ}{\mathbb{Q}}
\newcommand{\tr}{\mathrm{tr}}
\newcommand{\Rt}{R^{\mathrm{target}}}
\title{Gate 1: the $V_{1/2}$ heat-kernel coefficient $a_1$ is a clean $\int R$
on the soldered metric $g=e\!\cdot\!e$ over $\h_3(\Oo)$}
\date{2026-06-06 \quad (milestone v21.0, Phase 81; Gate 1 of the Sakharov ladder, run after Gate 0 SURVIVED)}
\begin{document}
\maketitle

\begin{abstract}
We verify, exactly over $\QQ$, that the Seeley--DeWitt coefficient $a_1$ obtained by
integrating out the $V_{1/2}=\mathbf{16}$ chiral Weyl fermions on the soldered Lorentzian
slice $g=e\!\cdot\!e$ of the exceptional Jordan algebra $\h_3(\Oo)$ is genuinely
proportional to the Ricci scalar $R[g]$ --- it is \emph{not} contaminated by a gauge
$\tr(F^2)$ term, by a wave-map target-curvature term, or by an octonionic-associator /
torsion correction. The single structural fact is that $a_1=\tr(E+\tfrac{1}{6}R)$ carries
\emph{only} the endomorphism $E$ and the scalar curvature $R\cdot\mathbb{1}$; every
$R_{\mu\nu}^2,\,R_{\mu\nu\rho\sigma}^2,\,R^2,\,\Box R,\,\tr F^2$ invariant lives one order
up, in $a_2$. The three candidate contamination channels are each ruled out: (i) the
internal gauge field strength drops out of $a_1$ because $\tr\gamma^{\mu\nu}=0$; (ii) the
wave-map target-curvature channel is absent because $V_{1/2}=\mathbf{16}$ of $\mathrm{Spin}(10)$
is a \emph{linear} spinor representation (flat target, $\Rt\equiv 0$); and (iii) the
soldering connection $\omega(e)$ is the torsion-free Levi--Civita connection (a
\emph{checked}, not assumed, fact: the torsion $2$-form vanishes identically over $\QQ$ on a
position-dependent tetrad), so the fluctuation operator is Laplace-type and $a_1$ has the
standard Gilkey form. The per-Dirac coefficient is $a_1=-\tfrac13 R$, the residual after
stripping the $R$-term is \emph{exactly} zero, and the read-off is $c_1=\STr/(4\pi)^2=
1/(6\pi^2)$, giving a candidate $\kappa^{-1}=c_1\Lambda_f^2$ and $G=3\pi/(8\Lambda_f^2)>0$
(consistent with the Gate-0 sign; $\Lambda_f$ is the Gate-3 scale, not pinned here). The
verdict is \textbf{PASS} ($a_1$ clean $\propto R$); Gate 1 does not kill, and the predicted
real death is deferred to Gate 2 (the v18 $16$-vs-$6$ tensor-support mismatch).
\end{abstract}

\section{The question and the structural theorem}\label{sec:setup}

Gate 0 of the v21.0 Sakharov ladder established (exactly over $\QQ$, verifier-hardened,
human-ratified) that the spin-weighted supertrace of the integrated-out $V_{1/2}$ sector is
$\STr=+8/3>0$, so the induced Newton constant is positive in the proper-time / cutoff
\emph{vacuum} scheme (Casimir-category; zero temperature; no thermodynamics). Gate 1 asks the
next question in the menu: is the induced curvature term \emph{clean}? That is, does the
Seeley--DeWitt expansion of the one-loop effective action produce a term proportional to the
spacetime Ricci scalar $R[g]$ --- the genuine Einstein--Hilbert structure --- or is it
polluted by an invariant that is not $\int R$, which would mean the route does not induce
\emph{gravity} at all but some other curvature functional?

The decisive lever is a theorem of Gilkey. For a Laplace-type operator $D^2=-(\nabla^2+E)$
on a vector bundle over a $d$-dimensional Riemannian manifold, the heat-kernel
(Seeley--DeWitt) coefficient at coincidence is
\begin{equation}\label{eq:gilkey-a1}
a_1(x,x)=\tr\!\left(E(x)+\tfrac{1}{6}R(x)\,\mathbb{1}\right),
\end{equation}
where $E$ is the bundle endomorphism, $R$ is the scalar curvature, $\mathbb{1}$ is the
bundle identity, and the trace is over the bundle indices (Vassilevich,
\textit{Phys.\ Rept.}\ \textbf{388} (2003) 279, \S4.3; Gilkey, \textit{Invariance Theory}).
The content of \eqref{eq:gilkey-a1} that drives this gate is \emph{structural}:
\begin{proposition}[$a_1$ carries only $E$ and $R$]\label{prop:a1-structure}
The coefficient $a_1$ contains \emph{only} the endomorphism $E$ and the scalar curvature
$R\cdot\mathbb{1}$. No quadratic curvature invariant
($R_{\mu\nu}^2$, $R_{\mu\nu\rho\sigma}^2$, $R^2$, $\Box R$) and no gauge invariant
$\tr F^2$ appears in $a_1$; all of these first appear one order up, in $a_2$ (Gilkey's
$a_4$).
\end{proposition}
Consequently, $a_1$ can fail to be a clean $\int R$ \emph{only} if the endomorphism $E$
itself carries a non-$R$ curvature invariant that survives the bundle trace with a nonzero
coefficient. There are exactly three candidate channels for such a contaminant, which we
treat in \S\S\ref{sec:gauge}--\ref{sec:torsion}, and dispatch in
\S\ref{sec:residual}--\S\ref{sec:verdict}.

\paragraph{Conventions (inherited, locked).} Metric signature mostly-minus
$(+,-,-,-)$; the soldered spacetime metric is $g=e\cdot e$ with $e=\pi_u(dE)$ on the
$\h_3(\Oo)$ Peirce slice (engine indices $\{1,2,3,10\}$), and the $V_{1/2}$ matter survivors
live on $\{11,18,19,26\}$. The $d=4$ Dirac spinor trace is $\tr\mathbb{1}=4=2^{d/2}$, and a
Weyl fermion is half the Dirac bundle. The squared Dirac operator is the Lichnerowicz form
$(i\gamma\!\cdot\!\nabla)^2=\nabla^2-R/4$, augmented by minimal gauge coupling (\S\ref{sec:gauge}).
All arithmetic is exact over $\QQ$ (and $\QQ\!\cdot\!\pi$ for the read-off); no float appears
on any decisive path.

\section{Channel 1: the gauge field strength drops out of $a_1$}\label{sec:gauge}

The minimally gauge-coupled squared Dirac operator is the Lichnerowicz--Weitzenb\"ock form
\begin{equation}\label{eq:lich-gauge}
(i\slashed{D})^2=-\nabla^2+\tfrac14 R+\tfrac12\,\gamma^{\mu\nu}F_{\mu\nu},
\qquad\Longrightarrow\qquad
E=-\tfrac14 R-\tfrac12\,\gamma^{\mu\nu}F_{\mu\nu},
\end{equation}
where $\nabla$ is the spacetime$+$gauge covariant derivative and $F_{\mu\nu}$ is the internal
field strength (here the $\mathfrak{so}(6)=\mathfrak{su}(4)$ curvature identified in
Phase~76 as the Berry / $\mathrm{Im}\,\mathrm{QGT}$ sector --- \emph{not} gravity).
Taking the bundle trace of $E+\tfrac16 R$ over the spinor$\otimes$gauge indices, the two
pieces of $E$ behave very differently:
\begin{align}
\tr\!\left(-\tfrac14 R\right)&=\tr(\mathbb{1})\cdot\left(-\tfrac14 R\right)=4\cdot\left(-\tfrac14 R\right)=-R\quad\text{(per Dirac)},\label{eq:lich-trace}\\
\tr\!\left(-\tfrac12\gamma^{\mu\nu}F_{\mu\nu}\right)&=-\tfrac12 F_{\mu\nu}\,\tr(\gamma^{\mu\nu})=0,\label{eq:F-trace}
\end{align}
because $\tr\gamma^{\mu\nu}=0$ identically: the antisymmetrized product of two distinct
gamma matrices is traceless in any dimension. Hence \emph{the gauge field strength drops out
of $a_1$}. The $\tr(F^2)$-type terms --- which would arise from $\tr(\gamma^{\mu\nu}
\gamma^{\rho\sigma})F_{\mu\nu}F_{\rho\sigma}\propto F^2$ --- appear only in $a_2$, through the
$\tfrac12\tr E^2$ piece, and merely renormalize the gauge kinetic term; they are irrelevant
to the induced Einstein--Hilbert coefficient. Combining \eqref{eq:lich-trace}--\eqref{eq:F-trace}
with the universal $+\tfrac16 R$ of \eqref{eq:gilkey-a1},
\begin{equation}\label{eq:a1-dirac}
a_1^{\text{Dirac}}=\tr\!\left(E+\tfrac16 R\right)
=4\left(-\tfrac14+\tfrac16\right)R
=4\cdot\left(-\tfrac{1}{12}\right)R
=-\tfrac13\,R,
\end{equation}
a clean multiple of $R$, with the per-Dirac-component coefficient $-\tfrac14+\tfrac16=-\tfrac{1}{12}$
and the per-Weyl coefficient $-\tfrac16$ (half the Dirac bundle). These reproduce the Gate-0
bundle values: with the Fermi statistics sign $s=-1$, the Weyl signed contribution is
$(-1)\cdot(-\tfrac16)=+\tfrac16$ and the Dirac signed contribution is $(-1)\cdot(-\tfrac13)=+\tfrac13$,
summing to $\STr=16\cdot(+\tfrac16)=+\tfrac83$.

\section{Channel 2: the wave-map target-curvature term is absent (linear rep)}\label{sec:wavemap}

The prompt writes the matter kinetic functional as $\langle DM,DM\rangle$ (``wave-map /
Dirac''). This is the one channel that requires genuine physical input, because a
\emph{nonlinear} sigma-model $\phi:\text{spacetime}\to M_{\mathrm{target}}$ into a
\emph{curved} target injects into the one-loop fluctuation endomorphism a term proportional
to the \emph{target} Riemann tensor contracted with background gradients (Friedan,
\textit{Ann.\ Phys.}\ \textbf{163} (1985) 318; Alvarez-Gaum\'e--Freedman--Mukhi,
\textit{Ann.\ Phys.}\ \textbf{134} (1981) 85):
\begin{equation}\label{eq:wavemap-contam}
E\;\supset\;-\,\Rt_{abcd}(\bar\phi)\,\partial_\mu\bar\phi^{\,b}\,\partial^\mu\bar\phi^{\,d}.
\end{equation}
This endomorphism is proportional to $\Rt(\partial\bar\phi)^2$, which is \emph{not}
proportional to the spacetime Ricci scalar $R[g]$. Traced into $a_1$ it would leave a
surviving non-$R$ invariant with generically nonzero coefficient, so $a_1$ would fail to be
a clean $\int R$ and the gate would die.

The decisive physical fact is that this channel is \emph{absent} for $V_{1/2}$:
\begin{lemma}[flat target]\label{lem:linear-rep}
$V_{1/2}=\mathbf{16}$ of $\mathrm{Spin}(10)$ is a \emph{linear} spinor representation --- one
Standard-Model generation of spin-$\tfrac12$ fermions (Paper~7; the $\h_3(\Oo)$ Peirce
half-eigenspace). The field is a section of a \emph{flat} fermionic bundle (the $16$-dimensional
spinor representation, equipped with the constant $\mathrm{Spin}(10)$-invariant bilinear)
twisted by the spacetime spin connection $\omega(e)$ and the internal $\mathrm{SU}(4)$ gauge
connection. The soldering $e=\pi_u(dE)$ acts \emph{linearly} on $V_{1/2}$; there is no curved
coset target. Hence the target Riemann tensor vanishes identically, $\Rt\equiv 0$, and the
contaminant \eqref{eq:wavemap-contam} is identically absent.
\end{lemma}
A linear representation is a vector space, hence flat: its ``target metric'' is the constant
invariant form. The word ``wave-map'' in the prompt is the generic name for the kinetic
functional $\langle DM,DM\rangle$; for a linear representation this is just the
(gauge$+$gravity)-twisted Dirac/Klein--Gordon kinetic term, with a flat target. Only a
genuine \emph{coset} sigma-model (e.g.\ an $E_6/(\cdots)$ scalar manifold) would inject target
curvature, and $V_{1/2}$ is not of that type. We verify computationally that the detection
machinery is \emph{not} a blind spot: feeding a synthetic curved-target term into $a_1$ produces
a nonzero $\partial a_1/\partial \Rt=-4\neq0$, confirming that the absence reported above is a
genuine finding rather than an undetectable case (the residual test of \S\ref{sec:residual}).

\section{Channel 3: the operator is Laplace-type (torsion-free, checked)}\label{sec:torsion}

The structural theorem \eqref{eq:gilkey-a1} presupposes that the fluctuation operator is
\emph{Laplace-type}, $D^2=-(\nabla^2+E)$ with $\nabla$ metric-compatible. This is the one
assumption that must be \emph{checked} rather than assumed, because the underlying algebra is
non-associative: a leftover octonionic-associator term in the covariant derivative $D$ of
$\langle DM,DM\rangle$ would render the operator non-Laplace-type, and the heat-kernel
coefficient of a non-Laplace-type operator does not reduce to $\tr(E+\tfrac16 R)$. Equally, a
nonzero torsion would add a $\propto T^2$ correction.

Milestone v18 (Phase~77) established, exactly over $\QQ$, that the soldering connection
$\omega(e)$ built from the tetrad $e=\pi_u(dE)$ (with $g=e\cdot e$) is the \emph{torsion-free}
Levi--Civita connection, and that the second-Cartan-structure curvature $R[\omega]$ equals the
metric Levi--Civita Riemann tensor of $g$ on every sampled component. We re-confirm the
torsion-free property here on a genuinely position-dependent tetrad, using the same harness
routines: with the rational warped-Lorentzian frame $e^a{}_\mu(x)=\mathrm{diag}(1,1+x_0^2,1,1)$,
the closed-form Levi--Civita connection
\begin{equation}\label{eq:omega-formula}
\omega_\mu{}^{ab}=\tfrac12 e^{\nu a}\!\left(\partial_\mu e_\nu^{\,b}-\partial_\nu e_\mu^{\,b}\right)
-\tfrac12 e^{\nu b}\!\left(\partial_\mu e_\nu^{\,a}-\partial_\nu e_\mu^{\,a}\right)
-\tfrac12 e^{\rho a}e^{\sigma b}\!\left(\partial_\rho e_\sigma^{\,c}-\partial_\sigma e_\rho^{\,c}\right)e_{\mu c}
\end{equation}
is antisymmetric, $\omega^{ab}=-\omega^{ba}$, and the torsion $2$-form
\begin{equation}\label{eq:torsion}
\Theta^a{}_{\mu\nu}=\partial_\mu e^a_\nu-\partial_\nu e^a_\mu
+\omega_\mu{}^a{}_b\,e^b_\nu-\omega_\nu{}^a{}_b\,e^b_\mu
\end{equation}
vanishes \emph{identically} over $\QQ$. Because the tetrad carries genuine $x_0$-dependence,
this is a non-trivial test: a non-Levi--Civita connection would yield a nonzero $\Theta$.
Hence the matter covariant derivative $D$ in $\langle DM,DM\rangle$ is the standard
metric-compatible $\nabla$ (plus the gauge connection), the operator is Laplace-type, and
$a_1$ has the standard Gilkey form with no associator and no $T^2$ correction.

\begin{remark}[distinct from the v20 Einstein--Cartan torsion]
The torsion that vanishes here is that of the \emph{matter Laplacian's} connection, the
Levi--Civita $\omega(e)$. It is a different object from the \emph{geometric} Einstein--Cartan
torsion of $e$ sourced by the chiral spin current, which milestone v20 found to be nonzero.
The two refer to different connections; the v20 finding does not reintroduce a $T^2$ term in
the matter heat kernel.
\end{remark}

\section{The $a_1\propto R$ residual test (decisive, exact over $\QQ$)}\label{sec:residual}

Assembling \eqref{eq:a1-dirac} with the (vanishing) channels of \S\S\ref{sec:gauge}--\ref{sec:torsion},
the full symbolic $a_1$ (carrying the gauge magnitude $F$ and a placeholder target-curvature
$\Rt$ for the synthetic test) is
\begin{equation}\label{eq:a1-full}
a_1=\tr(E+\tfrac16 R)= -\tfrac13\,R+0\cdot F+0\cdot \Rt .
\end{equation}
Decomposing $a_1$ into its coefficients on the curvature monomials it \emph{could} carry, and
forming the residual after stripping the wanted $R$-term, we find exactly over $\QQ$:
\begin{equation}\label{eq:residual}
\frac{\partial a_1}{\partial R}=-\tfrac13,\qquad
\frac{\partial a_1}{\partial F}=0,\qquad
\frac{\partial a_1}{\partial \Rt}=0,\qquad
a_1-\Big(\tfrac{\partial a_1}{\partial R}\Big)R=0,
\end{equation}
and $a_1$ is first-order in $R$ (degree $\le 1$), so no $R^2$, $\Box R$, or $R_{\mu\nu}^2$ piece
is present (consistent with Proposition~\ref{prop:a1-structure}). The residual is
\emph{exactly} zero: the only curvature structure surviving the bundle trace in $a_1$ is
$R\cdot(\text{scalar rational})$. As a guard against a blind spot, the same decomposition
applied to a synthetically contaminated $a_1$ (with the curved-target term of
\eqref{eq:wavemap-contam} switched on) returns $\partial a_1/\partial\Rt=-4\neq0$ and flags a
nonzero non-$R$ residual --- the test fires on contamination.

\section{Read-off of $c_1$ and the induced coupling}\label{sec:c1}

With $\STr=+8/3$ from Gate 0 and the explicit $(4\pi)^{-2}$ prefactor of the
quadratically-divergent $a_1$ term in $W=\pm\tfrac12\Tr\log D^2$, the induced Newton coupling is
\begin{equation}\label{eq:c1}
\frac{1}{16\pi G}=c_1\,\Lambda_f^2,\qquad
c_1=\frac{\STr}{(4\pi)^2}=\frac{8/3}{16\pi^2}=\frac{1}{6\pi^2},
\end{equation}
a rational multiple of $\pi^{-2}$, exact over $\QQ\!\cdot\!\pi$. The candidate induced
gravitational coupling, pinned up to the Gate-3 scale $\Lambda_f$ (the $\rho_J$ fixed-point
scale, out of scope here), is therefore
\begin{equation}\label{eq:kappa-G}
\kappa^{-1}\equiv\frac{1}{16\pi G}=\frac{1}{6\pi^2}\,\Lambda_f^2,
\qquad
G=\frac{1}{16\pi c_1\Lambda_f^2}=\frac{3\pi}{8\,\Lambda_f^2}>0,
\end{equation}
positive, consistent with the Gate-0 sign. The dimensionless content delivered by Gate 1 is
$c_1=1/(6\pi^2)$; the numerical magnitude awaits $\Lambda_f$ at Gate 3.

\section{Verdict}\label{sec:verdict}

The verdict is a deterministic function of three computed flags --- the non-$R$ residual, the
torsion, and the linearity of the representation --- not a hardwired label. The mapping is
\begin{equation}\label{eq:verdict}
\text{Gate 1}=
\begin{cases}
\text{DEAD (}a_1\text{ not clean }\propto R\text{)} & \text{if a non-}R\text{ invariant survives in }a_1,\\
\text{DEAD (non-Laplace-type)} & \text{if the connection has torsion / an associator,}\\
\text{DEAD (curved-target wave-map)} & \text{if }V_{1/2}\text{ is a curved sigma-model,}\\
\text{PASS (}a_1\text{ clean }\propto R\text{)} & \text{otherwise.}
\end{cases}
\end{equation}
The computed inputs are: non-$R$ residual $=0$ (\S\ref{sec:residual}), torsion $=0$
(\S\ref{sec:torsion}), linear representation / flat target (\S\ref{sec:wavemap}). All three
clean clauses hold, so
\begin{equation}\label{eq:pass}
\boxed{\ \text{Gate 1: PASS}\quad a_1=c_1\,\Lambda_f^2\,R=\frac{1}{6\pi^2}\,\Lambda_f^2\,R\quad(\text{clean }\propto R)\ }
\end{equation}
The three contamination channels are each independently demonstrated to be detectable
(synthetic injections flip the verdict to the corresponding DEAD branch), pre-empting a
hardwired-boolean false positive. Gate 1 therefore does \emph{not} kill the route: with a
positive sign (Gate 0) and a clean Einstein--Hilbert coefficient (Gate 1, here), the route
lives to Gate 2 --- the v18 tensor-support mismatch, the predicted real death (the induced
scalar $\kappa$ cannot repair the $16$-vs-$6$ component mismatch between $G[g]$ and $\kappa T$,
because a scalar coupling cannot change tensor rank).

\paragraph{Strength of the claim.} This is a PASS at true strength of \emph{the second of four
gates}, not a claim that ``Sakharov induced gravity works.'' Gate 1 certifies only that the
induced curvature term is a clean $\int R$; it makes no statement about whether the closure
$G[g]=\kappa T[M]$ holds (Gates 2--3). The $a_1=E+R/6$ structure is textbook (Gilkey;
Vassilevich), the three channel arguments are exact over $\QQ$ with the wave-map channel
ruled out by the linearity of the $\mathbf{16}$, and the one genuine risk (Laplace-type) is
checked, not assumed.

\section*{Sources}
\begin{itemize}
\item P.\ Gilkey, \textit{Invariance Theory, the Heat Equation, and the Atiyah--Singer
Index Theorem} --- the coefficient $a_2(x,x)=E+R/6$ (the $a_1$ structure: only $E$ and $R$),
and $a_4$ (where $R_{\mu\nu}^2,R_{\mu\nu\rho\sigma}^2,R^2,\Box R,F^2$ live).
\item D.\ V.\ Vassilevich, ``Heat kernel expansion: user's manual,'' \textit{Phys.\ Rept.}\
\textbf{388} (2003) 279 (hep-th/0306138), \S2 (Laplace-type requirement), \S4.3
($a_2=E+R/6$, the spinor/gauge endomorphisms, $\tr\gamma^{\mu\nu}=0$).
\item D.\ Friedan, ``Nonlinear models in $2+\varepsilon$ dimensions,'' \textit{Ann.\ Phys.}\
\textbf{163} (1985) 318 --- the sigma-model background-field heat kernel and the
target-Riemann endomorphism $\Rt(\partial\bar\phi)^2$.
\item L.\ Alvarez-Gaum\'e, D.\ Z.\ Freedman, S.\ Mukhi, \textit{Ann.\ Phys.}\ \textbf{134}
(1981) 85 --- the background-field method for sigma-model one-loop.
\item Frolov \& Fursaev, hep-th/9607104 --- the induced-$G$ weights cross-checked at Gate 0
($\STr=+8/3$).
\item \texttt{code/cartan\_phaseB\_curvature.py} (v18 Ph77) --- the torsion-free Levi--Civita
$\omega(e)$ and $R[\omega]=$ metric Riemann, exact over $\QQ$ (the \S\ref{sec:torsion} fact).
\item \texttt{code/bulk\_geometry\_verification.py} --- the exact-over-$\QQ$ $R[g=e\cdot e]\neq0$
anchor on the matter-bearing sample.
\end{itemize}

\end{document}
